\chapter{Observables from the ground state}
\label{chap:observables}

Let
\begin{equation}
\ket{\psi}=\sum_s \psi_s\ket{s},\qquad
\sum_s|\psi_s|^2=1.
\end{equation}
The observable routines normalize any finite nonzero input vector before
measurement.  This is convenient, but solver output should already have unit
norm; a large normalization correction would indicate an upstream problem.

\section{Diagonal observables}

Any occupation-only observable $O(s)$ is diagonal in the site basis:
\begin{equation}
\expect{O}=\sum_s |\psi_s|^2 O(s).
\end{equation}
This covers the total and local double occupancy,
\begin{align}
D &= \sum_i\expect{n_{i\uparrow}n_{i\downarrow}},\\
D/L &= \text{double occupancy per site},
\end{align}
as well as local charge and magnetization,
\begin{align}
\expect{n_i} &= \expect{n_{i\uparrow}+n_{i\downarrow}},\\
\expect{m_i} &= \expect{n_{i\uparrow}-n_{i\downarrow}}.
\end{align}
The project reports $m_i=2S_i^z$ for local magnetization and uses
\begin{equation}
S_i^z=\frac{n_{i\uparrow}-n_{i\downarrow}}{2}
\end{equation}
inside spin correlations.

The longitudinal spin and connected charge correlations are
\begin{align}
C^{zz}_{ij} &= \expect{S_i^zS_j^z},\\
C^n_{ij} &= \expect{n_in_j}-\expect{n_i}\expect{n_j}.
\end{align}
For fixed total particle number, summing $C^n_{ij}$ over both sites gives zero;
this is a useful sum rule.

\section{The one-body density matrix}

The spin-resolved one-body density matrix is not diagonal in the site basis:
\begin{equation}
\gamma^{(\sigma)}_{ij}
=\expect{c_{i\sigma}^\dagger c_{j\sigma}}.
\label{eq:obdm}
\end{equation}
For every source basis state, apply $c_{i\sigma}^\dagger c_{j\sigma}$ using the
same sign routine as the Hamiltonian, find the target basis state, and
accumulate
\begin{equation}
(\psi_{s'})^*\,\mathrm{sign}(s'\leftarrow s)\,\psi_s.
\end{equation}
The basic checks are
\begin{equation}
\gamma^{(\sigma)}=\left(\gamma^{(\sigma)}\right)^\dagger,
\qquad
\operatorname{Tr}\gamma^{(\sigma)}=N_\sigma.
\end{equation}
The eigenvectors of $\gamma$ are natural orbitals, but using them as a basis
for the interacting Hamiltonian requires the transformed interaction discussed
in Chapter~\ref{chap:generalhopping}.

\section{Fourier diagnostics}

On momenta $k_m=2\pi m/L$, the convention for the momentum distribution is
\begin{equation}
n_\sigma(k)=\frac{1}{L}\sum_{ij}e^{ik(i-j)}
\gamma^{(\sigma)}_{ij}.
\end{equation}
It obeys $\sum_k n_\sigma(k)=N_\sigma$.  The charge and longitudinal-spin
structure factors use the same transform:
\begin{align}
N(q)&=\frac{1}{L}\sum_{ij}e^{iq(i-j)}C^n_{ij},\\
S^{zz}(q)&=\frac{1}{L}\sum_{ij}e^{iq(i-j)}C^{zz}_{ij}.
\end{align}
The charge factor is connected by definition.  The spin factor is unconnected
by default and has an explicit connected option.

\lstinputlisting[caption={Measuring site, one-body, and Fourier observables.},label={lst:observables}]{code/step06_observables.py}

\paragraph{Checkpoint.}
Verify particle-number sums, Hermiticity and trace of the one-body density
matrix, reflection symmetry on an open chain, vanishing connected charge
$q=0$ weight, and nonnegative structure factors up to roundoff.
